% !TeX root = ../Book1.tex
\begin{figure}[htbp]
	\centering
	\begin{tikzpicture}[x=0.9cm,y=0.85cm,>=stealth,font=\small]
		\fill[GeometryBlue!5] (0,0) rectangle (3.2,2.4);
		\draw[GeometryGray] (0,0) rectangle (3.2,2.4);
		\foreach \z in {0.45,1.2,1.95}{\draw[GeometryBlue,thick,->] (0.2,\z)--(2.95,\z);}
		\draw[orange!80!black,thick] (0.9,0.12)--(0.9,2.28);
		\node[above] at (1.6,2.55) {时间与截面参数};
		\node[below] at (0.9,-0.12) {\(s=0\)};
		\node[below] at (2.75,-0.12) {\(s\)};
		\draw[->,thick] (3.7,1.2)--(5.2,1.2) node[midway,above] {\(F(s,z)\)};
		\begin{scope}[xshift=5.7cm]
			\draw[orange!80!black,thick] (0.65,0.05) .. controls (0.2,0.8) and (0.9,1.6) .. (0.6,2.45);
			\draw[GeometryBlue,thick,->] (0,0.35) .. controls (1,0.25) and (1.8,0.1) .. (3.15,0.9);
			\draw[GeometryBlue,thick,->] (0,1.1) .. controls (1,1.05) and (2,1.3) .. (3.15,1.65);
			\draw[GeometryBlue,thick,->] (0,1.95) .. controls (1,1.75) and (2,2.4) .. (3.15,2.25);
			\fill (0.54,1.1) circle (1.8pt);
			\node[below right] at (0.54,1.1) {\(p\)};
			\node[above] at (1.6,2.55) {流形中的局部轨道};
			\node[below] at (0.6,-0.12) {截面 \(\sigma\)};
		\end{scope}
	\end{tikzpicture}
	\caption[流盒与横截面]{流盒与横截面。橙线给出初始截面；蓝线是从截面出发的局部积分曲线，箭头表示正时间。映射 \(F(s,z)=\Phi_s(\sigma(z))\) 把左侧的时间平移变成右侧沿向量场的运动，图形仅示意局部构造。}
	\label{g1:fig:flow-box}
\end{figure}
